% 36-31 ②(36쪽) — 보기 그래프: x<a 는 증가하는 직선(x=a 에서 채운 점), x>a 는 가로 반직선(x=a 에서 빈 점)으로 뛰어 끊긴 그래프.
% 보기 자체가 그림이라 TikZ 로 다시 그림. 라벨은 원문 그대로 y=f(x), O, a, x, y 만.
\begin{tikzpicture}[scale=0.62, line width=0.5pt, baseline=(current bounding box.center)]
  \path (-2.3,-1.6) rectangle (3.1,2.7);
  \draw[->, >=stealth, line width=0.7pt] (-2.1,0) -- (2.85,0) node[below=1pt, font=\scriptsize] {$x$};
  \draw[->, >=stealth, line width=0.7pt] (0,-1.15) -- (0,2.35) node[left=1pt, font=\scriptsize] {$y$};
  \draw[line width=0.6pt] (-2.0,-1.145) -- (0.75,0);
  \draw[line width=0.6pt] (0.75,1.15) -- (2.7,1.15);
  \draw[densely dotted, line width=0.4pt] (0.75,0) -- (0.75,1.15);
  \fill (0.75,0) circle (1.8pt);
  \draw[fill=white] (0.75,1.15) circle (1.8pt);
  \node[anchor=south, font=\scriptsize] at (1.75,1.30) {$y=f(x)$};
  \node[below left=-2pt and -2pt, font=\scriptsize] at (0,0) {$\pt{O}$};
  \node[below=1pt, font=\scriptsize] at (0.75,0) {$a$};
\end{tikzpicture}
